\section{Methodology}
\label{sec:methodology}

This section formalizes the ink detection pipeline, covering data representation, Z-layer sampling, training procedure, and evaluation metrics.

\subsection{Data Representation and Z-Layer Geometry}
\label{sec:data_representation}

Each scroll segment is produced by the segmentation step of the virtual unwrapping pipeline~\cite{seales2023educelab}. A segment represents a locally flattened patch of papyrus surface, described by a triangulated mesh in 3D CT space. For each mesh vertex, the segmentation tool samples the CT volume at 65 uniformly spaced depth offsets perpendicular to the estimated surface normal, producing a stack of 2D grayscale images:

\begin{equation}
\label{eq:zlayer_stack}
\mathbf{S} = \{I_z \in \mathbb{R}^{H \times W} \mid z = 0, 1, \ldots, 64\}
\end{equation}

\noindent where $H$ and $W$ are the spatial dimensions of the segment (varying per segment, typically $500$--$4000$ pixels), and each layer $I_z$ represents a depth offset of $z \times 7.91$~\textmu m from an arbitrary reference point. Layer $z = 32$ is conventionally aligned with the estimated papyrus surface center.

The physical structure of the papyrus at the scan site spans approximately 19 layers (${\sim}150$~\textmu m), occupying layers 23-41 in the 65-layer stack. Ink, applied to the papyrus surface in antiquity, is expected to produce a density signal in layers 30-35, corresponding to the surface and the first few micrometers of ink penetration into the papyrus fibers.

\subsubsection{Z-Layer Sampling}

Given the full 65-layer stack, the Z-layer sampling function selects a contiguous window of $N$ layers starting at index $s$:

\begin{equation}
\label{eq:zlayer_window}
\mathbf{X} = [I_s, I_{s+1}, \ldots, I_{s+N-1}] \in \mathbb{R}^{N \times H \times W}
\end{equation}

The default configuration uses $s = 30$ and $N = 5$, selecting layers 30-34. The Grand Prize solution uses $s = 17$ and $N = 26$, selecting layers 17-42. The choice of $(s, N)$ is the central independent variable of this study.

\subsubsection{Intensity Preprocessing}

Raw CT values are clipped and normalized as $\hat{I}_z(x, y) = \min(I_z(x, y),\,\tau)/\tau$ with $\tau = 200$, mapping all pixels to $[0, 1]$ and suppressing metal artifacts above the ink-relevant density range.

\subsubsection{Reversed Segments}

Eight segments have inverted Z-ordering (layer 0 faces the scroll exterior). These are identified and their layer stacks are reversed prior to processing to ensure consistent depth semantics across all segments.

\subsection{Tiling Strategy}
\label{sec:tiling}

Full segments are too large for direct model processing. We decompose each segment into overlapping square tiles:

\begin{equation}
\mathbf{t}_{r,c} = \mathbf{X}[:, r:r+P, c:c+P] \in \mathbb{R}^{N \times P \times P}
\end{equation}

\noindent where $P = 64$ is the tile size and $(r, c)$ are the top-left coordinates sampled at stride $\delta$. During training, $\delta = 32$ (50\% overlap); during inference, $\delta$ is configurable (typically 21--32). Segment boundaries are padded to the nearest multiple of 256 pixels using zero-padding prior to tiling.

\subsection{Training Procedure}
\label{sec:training}

\subsubsection{Dataset and Splits}

Training data consists of tiles extracted from four labeled segments of Scroll~1: 20230504231922, 20230504225948, 20230504171956, and 20230504125349. Segment 20230504093154 is reserved for validation. Segment 20230504223647 serves as the held-out test set, never seen during training or hyperparameter selection.

Ground truth labels are pseudo-labels generated by the Grand Prize TimeSformer model, binarized at threshold 0.4.

\subsubsection{Pseudo-Label Generation}
\label{sec:pseudo_labels}

Manual ink annotation requires papyrology expertise and is prohibitively expensive at scale. We adopt a pseudo-label approach: the pre-trained Grand Prize TimeSformer, trained on the official Vesuvius Challenge labels, generates predictions on all segments using its canonical 26-layer configuration. These predictions are binarized at threshold 0.4 to produce binary ink labels, and stored as soft probability maps for potential use in distillation.

This approach is valid because: (a) the GP model achieved verified ink detection accuracy sufficient to win the Grand Prize, (b) the pseudo-labels are deterministic (fixed model, fixed weights), and (c) the approach functions as knowledge distillation, where the GP model is the teacher and our ablation models are students.

\subsubsection{Loss Function}

We employ a composite loss combining Dice loss and Soft Binary Cross-Entropy (SoftBCE) with label smoothing:

\begin{equation}
\label{eq:loss}
\mathcal{L} = \alpha \cdot \mathcal{L}_{\text{Dice}} + \beta \cdot \mathcal{L}_{\text{SoftBCE}}
\end{equation}

\noindent with $\alpha = \beta = 0.5$.

The Dice loss directly optimizes the overlap between prediction and target:

\begin{equation}
\mathcal{L}_{\text{Dice}} = 1 - \frac{2 \sum_i p_i g_i + \epsilon}{\sum_i p_i + \sum_i g_i + \epsilon}
\end{equation}

\noindent where $p_i = \sigma(\hat{y}_i)$ is the predicted probability (after sigmoid), $g_i$ is the ground truth label, and $\epsilon$ is a smoothing constant for numerical stability.

The SoftBCE loss applies label smoothing with factor $\gamma = 0.25$:

\begin{equation}
\mathcal{L}_{\text{SoftBCE}} = -\frac{1}{|T|}\sum_i \left[\tilde{g}_i \log p_i + (1 - \tilde{g}_i) \log(1 - p_i)\right]
\end{equation}

\noindent where $\tilde{g}_i = g_i(1 - \gamma) + 0.5\gamma$ is the smoothed label.

This combination provides the complementary benefits of Dice loss (robust to class imbalance, directly optimizes overlap) and BCE (pixel-wise gradient signal, smooth optimization landscape). Label smoothing prevents the model from becoming overconfident in regions where the pseudo-labels may be imprecise.

\subsubsection{Optimization}

Models are trained using the AdamW optimizer~\cite{loshchilov2019adamw} with decoupled weight decay:

\begin{itemize}
    \item Learning rate: $\eta = 3 \times 10^{-4}$
    \item Weight decay: $\lambda = 10^{-5}$
    \item Gradient clipping: $\|\nabla\|_{\max} = 5.0$
\end{itemize}

The learning rate follows a warmup-cosine schedule:

\begin{equation}
\eta(t) = \begin{cases}
\eta_{\max} \cdot \frac{t}{T_w} & \text{if } t \leq T_w \\[3pt]
\frac{\eta_{\max}}{2} \left(1 + \cos\!\left(\frac{\pi (t - T_w)}{T - T_w}\right)\right) & \text{if } t > T_w
\end{cases}
\end{equation}

\noindent where $T_w = 3$ is the warmup duration in epochs and $T$ is the total number of training epochs.

\subsubsection{Data Augmentation}

Standard spatial augmentations (rotation, flips, scaling, elastic deformation, brightness/contrast, Gaussian noise) are applied identically across all Z-layers within a tile. Three depth-specific Z-axis augmentations address physical variability: Z-flip ($p = 0.5$) simulates inverted scanning orientations, Z-shift ($\Delta z \in \{-2,\ldots,+2\}$, $p = 0.3$) simulates surface estimation uncertainty, and layer dropout ($p = 0.1$) trains robustness to missing depth data.

\subsection{Evaluation Metrics}
\label{sec:metrics}

Following the Vesuvius Challenge convention~\cite{vesuvius2023challenge}, the primary metric is the $F_{0.5}$ score, which weights precision twice as heavily as recall:

\begin{equation}
\label{eq:f05}
F_{0.5} = \frac{(1 + 0.5^2)\,\text{Precision} \cdot \text{Recall}}{0.5^2\,\text{Precision} + \text{Recall}} = \frac{1.25\,P \cdot R}{0.25\,P + R}
\end{equation}

\noindent This emphasis on precision reflects the application context: false ink detections mislead downstream historical interpretation more than missed ink regions. Predictions are binarized at threshold $\theta = 0.4$ before scoring.

Secondary metrics include the Dice coefficient (symmetric overlap without precision weighting), per-pixel precision and recall, mean prediction intensity (a proxy for prediction confidence when ground truth is unavailable), and coverage at $\theta = 0.4$ (the fraction of pixels whose predicted probability exceeds the binarization threshold).
